\dang{BT đề cương}
\begin{ex}
Một chất điểm dao động điều hoà có chu kì $T=1\text{ s}$. Tần số góc $\omega$ của dao động là
\choice
{$\pi \text{ (rad/s)}$}
{\True $2\pi \text{ (rad/s)}$}
{$1 \text{ (rad/s)}$}
{$2 \text{ (rad/s)}$}
\loigiai{Ta có công thức liên hệ giữa tần số góc và chu kì: $\omega = \dfrac{2\pi}{T}$.
Thay số: $\omega = \dfrac{2\pi}{1} = 2\pi \text{ (rad/s)}$.}
\end{ex}
\dang{BT đề cương}
\begin{ex}Một chất điểm dao động điều hòa với tần số góc $\omega=10\pi \text{ (rad/s)}$. Tần số của dao động là
\choice
{ \True $5 \text{ Hz}$}
{ $10 \text{ Hz}$}
{ $20 \text{ Hz}$}
{ $5\pi \text{ Hz}$}
\loigiai{
Ta có công thức liên hệ giữa tần số và tần số góc là $f = \dfrac{\omega}{2\pi}$.
Thay số vào công thức, ta được $f = \dfrac{10\pi}{2\pi} = 5 \text{ Hz}$.
Vậy tần số của dao động là $5 \text{ Hz}$.
}
\end{ex}
\dang{BT đề cương}
\begin{ex}
Một chất điểm dao động điều hoà. Trong thời gian 1 phút, vật thực hiện được 30 dao động. Chu kì dao động của chất điểm là
\choice
{\True $2 \text{ s}$}
{$30 \text{ s}$}
{ $0,5 \text{ s}$}
{$1 \text{ s}$}
\loigiai{Thời gian thực hiện dao động là $\Delta t = 1 \text{ phút} = 60 \text{ s}$.
Số dao động thực hiện được là $N = 30$.
Chu kì dao động được tính bằng công thức $T = \dfrac{\Delta t}{N}$.
Thay số: $T = \dfrac{60}{30} = 2 \text{ s}$.

\end{ex}
\dang{BT đề cương}
\begin{ex}
Một chất điểm dao động điều hoà có phương trình li độ theo thời gian là $x=5\sqrt{3}\cos \left(10\pi t+\dfrac{\pi}{3}\right)\text{ (cm)}$. Tần số của dao động là
\choice
{ $10 \text{ Hz}$}
{$20 \text{ Hz}$}
{$10\pi \text{ Hz}$}
{\True $5 \text{ Hz}$}
\loigiai{Từ phương trình dao động điều hoà $x=A\cos(\omega t+\varphi)$, ta có tần số góc $\omega = 10\pi \text{ rad/s}$.
Tần số của dao động được tính bằng công thức $f = \dfrac{\omega}{2\pi}$.
Thay số: $f = \dfrac{10\pi}{2\pi} = 5 \text{ Hz}$.}
\end{ex}
\dang{BT đề cương}
\begin{ex}
Một chất điểm dao động điều hoà có phương trình li độ theo thời gian là $x=6\cos \left(4\pi t+\dfrac{\pi}{3}\right)\text{ (cm)}$. Chu kì của dao động là
\choice
{$4 \text{ s}$}
{$2 \text{ s}$}
{$0,25 \text{ cm}$}
{\True $0,5 \text{ s}$}
\loigiai{Từ phương trình dao động điều hoà $x=A\cos(\omega t+\varphi)$, ta có tần số góc $\omega = 4\pi \text{ rad/s}$.
Chu kì của dao động được tính bằng công thức $T = \dfrac{2\pi}{\omega}$.
Thay số: $T = \dfrac{2\pi}{4\pi} = 0,5 \text{ s}$.}
\end{ex}
\dang{BT đề cương}
\begin{ex}
Một chất điểm dao động điều hoà có phương trình li độ theo thời gian là $x=10\cos \left(\dfrac{\pi}{3}t+\dfrac{\pi}{2}\right)\text{ (cm)}$.
Tại thời điểm $t$ vật có li độ $6 \text{ cm}$ và đang hướng về vị trí cân bằng. Sau $9 \text{ s}$ kể từ thời điểm $t$ thì vật đi qua li độ
\choice
{$3 \text{ cm}$ đang hướng về vị trí cân bằng}
{$-3 \text{ cm}$ đang hướng về vị trí biên}
{$6 \text{ cm}$ đang hướng về vị trí biên.}
{\True $-6 \text{ cm}$ đang hướng về vị trí cân bằng.}
\loigiai{Từ phương trình dao động $x=10\cos \left(\dfrac{\pi}{3}t+\dfrac{\pi}{2}\right)\text{ (cm)}$, ta có biên độ $A=10 \text{ cm}$ và tần số góc $\omega = \dfrac{\pi}{3} \text{ rad/s}$.
Chu kì của dao động là $T = \dfrac{2\pi}{\omega} = \dfrac{2\pi}{\pi/3} = 6 \text{ s}$.
Tại thời điểm $t$, vật có li độ $x=6 \text{ cm}$ và đang hướng về vị trí cân bằng. Điều này có nghĩa là vật đang chuyển động theo chiều âm (vì $x>0$ và hướng về VTCB).
Sau $9 \text{ s}$ kể từ thời điểm $t$, tức là tại thời điểm $t' = t + 9$.
Ta có $9 \text{ s} = 1,5 T$.
Sau $1T$, vật trở về trạng thái ban đầu (li độ và chiều chuyển động).
Sau $0,5T$, vật sẽ ở vị trí đối xứng với vị trí ban đầu qua gốc tọa độ và có chiều chuyển động ngược lại.
Vậy sau $1,5T$, vật sẽ ở vị trí đối xứng với vị trí ban đầu qua gốc tọa độ và có chiều chuyển động ngược lại so với trạng thái sau $1T$.
Tại thời điểm $t$, $x=6 \text{ cm}$ và đang hướng về VTCB (chiều âm).
Sau $1T$, vật ở $x=6 \text{ cm}$ và đang hướng về VTCB (chiều âm).
Sau thêm $0,5T$ (tức là tổng cộng $1,5T$), vật sẽ ở li độ $x' = -x = -6 \text{ cm}$ và đang hướng về VTCB (chiều dương).
Vậy sau $9 \text{ s}$, vật đi qua li độ $-6 \text{ cm}$ và đang hướng về vị trí cân bằng.}
\end{ex}
